\section{Сравнение с существующими аналогами}

Разработанная система прогнозирует волатильность с помощью альтернативных
источников данных. Ниже рассмотрены наиболее близкие аналоги~---
коммерческие платформы, открытые библиотеки и академические системы.

\subsection{Коммерческие платформы}

\textbf{Bloomberg Terminal / Refinitiv Eikon.}
Профессиональные терминалы предоставляют широкий набор инструментов
для анализа волатильности: реализованную волатильность, подразумеваемую волатильность,
исторические ряды VIX-подобных индексов.
Тем не менее эти платформы ориентированы на западные рынки,
не имеют готовых sentiment-моделей для русскоязычного контента,
а доступ к данным по инструментам MOEX после 2022~г.\ существенно ограничен.
Стоимость подписки ($>$\$20\,000 в год) исключает использование в академических проектах.

\textbf{MOEX Analytics.}
Официальный аналитический портал Московской биржи предоставляет
исторические данные OHLCV и простые индикаторы волатильности
(скользящее стандартное отклонение).
Инструменты машинного обучения, альтернативные источники данных
и функционал прогнозирования в платформе отсутствуют.

\subsection{Открытые библиотеки}

\textbf{arch (Python).}
Библиотека~\texttt{arch}~\cite{Sheppard2017} реализует широкий спектр
GARCH-семейства моделей (GARCH, EGARCH, GJR-GARCH и др.)
и является стандартом для параметрического моделирования волатильности.
Однако библиотека не поддерживает HAR-RV модели, не имеет встроенного
сбора данных и не работает с альтернативными источниками.
Настоящая система использует HAR-RV как бенчмарк,
а не GARCH, что обусловлено недоступностью высокочастотных данных.

\textbf{FinRL (AI4Finance).}
Фреймворк для обучения с подкреплением в финансах~\cite{Liu2021}
предоставляет среды для торговли акциями с базовыми техническими признаками.
FinRL не ориентирован на задачу прогнозирования волатильности,
не поддерживает sentiment-признаки из внешних API
и не имеет модулей для работы с MOEX.

\textbf{Qlib (Microsoft).}
Исследовательская платформа для количественных финансов~\cite{Yang2020}
включает обширный набор ML-моделей для прогнозирования доходностей.
Поддержка данных ограничена американскими и китайскими рынками;
модуль прогнозирования волатильности отсутствует;
русскоязычный sentiment не поддерживается.

\subsection{Академические системы}

В работе~\cite{Halouskova2025} предложена система прогнозирования
волатильности акций США на основе вектора макроэкономического внимания.
Настоящий проект является адаптацией этой методологии
для российского рынка драгоценных металлов:
заменены источники внимания (американские макроанонсы~→
решения ЦБ РФ, русскоязычные Telegram-каналы, GDELT-тональность)
и рыночные данные (S\&P~500~→ MOEX GLDRUB\_TOM / SLVRUB\_TOM).

\subsection{Сводная таблица сравнения}

\begin{table}[H]
\centering
\caption{Сравнение разработанной системы с аналогами}
\label{tab:analogs}
\small
\begin{tabular}{p{3.4cm}ccccc}
\hline
Характеристика
  & \rotatebox{70}{Bloomberg}
  & \rotatebox{70}{arch}
  & \rotatebox{70}{FinRL}
  & \rotatebox{70}{Qlib}
  & \rotatebox{70}{\textbf{Наша система}} \\
\hline
Данные MOEX              & $\pm$ & ---  & --- & --- & \textbf{+} \\
HAR-RV модель            & ---   & ---  & --- & --- & \textbf{+} \\
XGBoost / kNN            & ---   & ---  & +   & +   & \textbf{+} \\
Русскоязычный sentiment  & ---   & ---  & --- & --- & \textbf{+} \\
Telegram-парсер          & ---   & ---  & --- & --- & \textbf{+} \\
GDELT-тональность        & $\pm$ & ---  & --- & --- & \textbf{+} \\
Google Trends attention  & ---   & ---  & --- & --- & \textbf{+} \\
Макро ЦБ РФ              & ---   & ---  & --- & --- & \textbf{+} \\
Инкрементальное обновление& +    & ---  & --- & +   & \textbf{+} \\
Веб-интерфейс            & +     & ---  & --- & $\pm$ & \textbf{+} \\
Открытый исходный код    & ---   & +    & +   & +   & \textbf{+} \\
Бесплатный доступ        & ---   & +    & +   & +   & \textbf{+} \\
\hline
\end{tabular}

\medskip
\raggedright
\footnotesize
\textit{Обозначения:} «+»~--- поддерживается; «---»~--- отсутствует;
«$\pm$»~--- частичная поддержка.
\end{table}

Видно, что среди рассмотренных решений нет ни одного, которое
одновременно работает с данными MOEX, русскоязычным sentiment
и моделями прогнозирования волатильности (HAR-RV, бустинг, kNN).
Именно это сочетание и отличает разработанную систему от существующих
аналогов для российского рынка драгоценных металлов.
